% This is the chapter 'Physics-Introductory' to the "Hodgkin-Huxley neuron model V2.0 book"
% Last edited 2026.06.05
\section{Bevezetés\label{sec:Physics-Introduction}}
Az élő anyag teljes körű tudományos leírása minden \textit{diszciplináris} próbálkozásnak ellenáll: az élet létezése sok szempontból ellentmond a tudomány (rosszul alkalmazott) törvényeinek, amelyeket az élettelen természet objektumain teszteltünk. Saját törvényei vannak, amelyek ugyan hasonlítanak az élettelen természettudomány jól ismert törvényeire, de nem azonosak velük. Régóta ismert probléma, amit E.Schrödinger, Nobel-díjas fizikus, úgy foglalt össze nyolc évtizeddel ezelőtt, hogy az élő anyagot nem lehet leírni a fizika \textit{rendes} (ordinary) törvényeivel. Hozzátette azonban, hogy semmiféle új erő vagy kölcsönhatás nem lép fel az élő anyagban, csupán új összefüggéseket (más módszerű megértést) kell találni, és ezután a felismert új törvényszerűségek éppúgy illeszkedni fognak a tudomány szövetébe, mint a régiek. Hasonló véleményt mondott R.P.Feynman, szintén Nobel-díjas fizikus, is, amikor “kirándulásokat” javasolt más tudományterületekre, mivel az igazán érdekes jelenségek a rendes határokon túl, a “senki földjén” történnek.
Azaz, a (V2.0) tárgyalását megalapozandó, fizikát is tárgyalnunk kell,
bizonyos fogalmait, közelítéseit, törvényeit újraértelmezve, de mindig tiszteletben tartva az ún. első elveit. Ez helyenként nem-hagyományos gondolkodást és esetenként megszokott következtetések újra gondolását igényli.



\ifx\magyar\undefined
 Knowledge of the fundamental (physical) functioning of the biological neuron is essential for understanding its electrical processes and the excitation process, information processing based on signal transmission between neurons, and the low-level functioning of the brain. Although the fact that neurons show electrically measurable changes (i.e., belonging to the subject of electrical science) has been known since the beginning of instrumental studies, it was soon demonstrated that changes falling within the subject of thermodynamics and mechanics (i.e., another, separate discipline of physics) can also be measured. In the traditional view of physics, each discipline interprets only a certain range of quantities and can only derive relationships and laws between them. Nature, of course, does not care about the disciplinary nature of science, so science is essentially puzzled by the “miracle” that during the operation of a neuron interpreted as a purely electrical model, quantities falling within the scope of other disciplines (temperature, pressure) also change. The logic of things follows that then one must try to describe the experiences with the tools and laws of each discipline. In this way, it is possible to perform a “complete measurement”,
 \index{complete measurement}
i.e. to experimentally measure all quantities that are related to the operation of the neuron, but no connection is established between the disciplinary laws and quantities.
\else
A biológiai neuron alapvető (fizikai folyamatokon alapuló) működésének ismerete elengedhetetlen az elektromos folyamatoknak és a gerjesztés folyamatának, a neuronok közötti jelátvitelen alapuló információfeldolgozásnak, és az agy alacsony szintű működésének megértéséhez. Bár az a tény, hogy a neuronok elektromosan mérhető (azaz az elektromosságtan témakörébe tartozó) változásokat mutatnak, a műszeres vizsgálatok kezdete óta ismert, nem sokkal később azt is kimutatták, hogy a termodinamika és a mechanika (azaz a fizika másik, különálló diszciplinája) témakörébe eső változások is mérhetők. A fizika hagyományos szemléletében egy-egy diszciplina csupán bizonyos  mennyiségek körét értelmezi, csak azok közötti összefüggéseket és törvényeket tud származtatni. A természet természetesen nem törődik a tudomány diszciplinaritásával, így a tudomány lényegében értetlenül áll ama “csoda” előtt, hogy a tisztán elektromos modellként értelmezett neuron működése során más diszciplina tárgykörébe eső mennyiségek (hőmérséklet, nyomás) is megváltoznak. A dolgok logikájából következik, hogy akkor mindegyik diszciplina eszközeivel és törvényeivel kell megpróbálni leírni a tapasztalatokat. Ilyen módon lehet ugyan “teljes mérést” végezni, 
\index{teljes mérés}
azaz kísérletileg minden olyan mennyiséget mérni, ami a neuron működéséhez kapcsolódik, de a diszciplináris törvények és mennyiségek között nem jön létre kapcsolat.
\fi

\ifx\magyar\undefined
Although it is obvious that all phenomena are related to the same physical process, and the quantities are clearly correlated, a qualitative step forward is needed that does not simply put together the disciplinary laws developed for inanimate science, but rather adapts them to the conditions of living matter and creates the necessary connections between them. The infinitely distant, infinitely small change-causing concepts of physics, and disciplinary separation, are certainly not applicable. The amount of electric charges changing their location during operation is of the same order of magnitude as the amount of free charges on the membrane surface, so we can only use Lotka-Volterra type differential equations
\index{Lotka-Volterra type differential equation}
to describe electrical changes (the usual use of “equivalent current circuits” is incorrect and misleading), and furthermore, between the closed biological boundary walls of neurons, the interaction of ions on each other creates a counterforce on the wall and the change in the counterforce creates a pressure change in the electrolyte filling the cell. The elastic lipid layer forming the neuron wall undergoes a size change due to the pressure change (although it has a constant modulus of elasticity: during operation, the neuron’s size of about 25 micrometers changes by only about 1 nanometer).
\else
Bár nyilvánvaló, hogy minden jelenség ugyanahhoz a fizikai folyamathoz kapcsolódik, és a mennyiségek jól láthatóan korrelálnak, olyan minőségi továbblépésre van szükség, ami nem egyszerűen egymás mellé teszi az élettelen tudományra kidolgozott diszciplináris törvényeket, hanem az élő anyag körülményeire átdolgozza azokat és megteremti közöttük a szükséges kapcsolódásokat. Biztosan nem alkalmazhatók a fizika végtelen távol levő, végtelen kis változást okozó fogalmai, és a diszciplináris elkülönülés. A működés során helyüket megváltoztató elektromos töltések mennyisége ugyanabba a nagyságrendbe esik, mint a már/még ott levő szabad töltések mennyisége, így az elektromos változások leírására csupán Lotka-Volterra típusú differenciálegyenleteket használhatunk
\index{Lotka-Volterra differenciálegyenlet} (a szokásos “ekvivalens áramörök” használata hibás és félrevezető), továbbá a neuronok zárt biológiai határfalai között az ionok egymásra gyakorolt hatása a falon ellenerőt kelt és az ellenerő változása a sejtet kitöltő elektrolitban nyomásváltozást hoz létre. A neuron falát alkotó rugalmas lipidréteg a nyomásváltozás hatására méretváltozást szenved (bár órási rugalmassági modulusa van: a működés során a neuron kb 25 mikrométeres mérete csupán kb 1 nanométernyit változik). 
\fi


The roles of time, space and matter are subjects of endless debates in science.
Considering finite interaction speeds is against using a "nice and classical physics"
with its nice mathematical formulas, but omitting the different speeds
misled and may mislead research in several fields.
Biology produces situations where the complexity of phenomena
and the needed carefulness meets the ones needed in cosmology.
The difference is that, in biology, the consequences of phenomena are immediate
and they can be studied experimentally.


To describe the related phenomena, we must scrutinize, case by case, which interactions are significant and
which interaction(s) can be omitted; instead of setting up ad-hoc models for living matter that contradict
each other and fundamental scientific principles when used outside their narrow range of validity. To provide their
correct physics-based description, we must understand the corresponding behavior of living material, including that it works with slow ion currents,
electrically active, non-isotropic, structured materials,
and consequently, its temporal behavior (the speeds of interactions)
matters. We must consider macroscopic and microscopic phenomena (continuous and discrete features) at the same time, different
science fields, and their interplay (or mixing). "\href{https://www.pnas.org/doi/10.1073/pnas.1705704114}{Living complex systems} in
particular create high-ordered functionalities by pairing up low-ordered
complementary processes, e.g., one process to build and the other
to correct".~\cite{BiologicalConservationLaw:2017}
\index{law of conservation in biology~\cite{BiologicalConservationLaw:2017}}
We must check the validity of our abstractions.


First of all, we introduce (and, in this aspect, correct the current common understanding)
that neurons' "output product"~\cite{BuzsakiCellAssemblies:2010}
\textit{is created and transmitted by thermoelectric processes instead of
	net electric phenomena or net thermodinamic}.
	Feynman was right, also in this point: "\textit{nature is not interested in our separations, and many of the interesting phenomena bridge the gaps between fields.}" ~\cite{FeynmanThinking:1980}
We derive the correct mathematical and physical handling processes,
not considering the disciplinary boundaries, pointing to where
the \hyperlink{Abstractions}{disciplinary omissions and approximations} are not valid,
deriving the required correct handling of the physical processes together with with the mathematics required to handle them.
We use the theory developed in section~\ref{sec:Physics-Thermodynamics},
where (among others) \hyperlink{NernstPlanckTimeDerivatives}{the time derivatives of the processes} are introduced, in this way enabling the
quantitative mathematical description of the neurons' abstract operation,
that serves as a basis for discussing quantitatively the neural computing
processes (as we discuss in section~\ref{sec:Physics-LawsOfMotion},
these derivatives are the laws of motion of ions) and their information handling. However, these corrections must be carried out
in disciplinarily separated chapters.  
\index{laws of motion}


As Schr\"odinger and Feynman implicitly suggested, we revisit the approximations that led to classical physics (where we derived the well-known 'ordinary' laws).
\index{ordinary laws}
We scrutinize the "construction" and "working" of living matter in a cross-disciplinary way, using non-ordinary approximations and abstractions.
Then, laws based on the same first principles in a different approximation can describe living matter.
We show that nature and its observable phenomena are complex; it is pointless to dispute
which of the scientific disciplines describes neuronal behavior. The non-disciplinary approach shows that the firm and fast electric interaction sufficiently well describes neuronal operation after introducing the equivalent "thermodynamic electrical potential".
\index{thermodynamic electrical potential}



